\section{Purpose: Lepton Masses from One Parameter}

\textbf{[K]} The lepton sector is where the theory performs most strongly
and is most sharply testable. From the single parameter $\xi = 4/30000$ and
a rational assignment per lepton, the three charged lepton masses follow, their
ratios, the Koide relation $Q=2/3$, and -- the sharpest point -- a
prediction of the $\tau$-mass. This document develops both: the \emph{ladder} that
generates the masses, and the \emph{circulant} that sharpens the Koide structure.

\section{The Yukawa Ladder}

\textbf{[K]} The basic mass formula is
\[
  m_i = r_i\,\xi^{p_i}\,v, \qquad v = 246{.}22\ \text{GeV (Higgs-VEV)},
\]
with rational prefactors $r_i$ and rational exponents $p_i$ per lepton:
\[
  \begin{array}{lccc}
    \text{Lepton} & r_i & p_i & m_i\ (\text{calculated}) \\[2pt]
    \text{Electron} & 4/3 & 3/2 & 0{.}505\ \text{MeV} \\
    \text{Muon}     & 16/5 & 1  & 105{.}1\ \text{MeV} \\
    \text{Tau}      & 25/9 & 2/3 & 1785\ \text{MeV}
  \end{array}
\]

\textbf{[K]} The exponents form a ladder with fixed steps:
\[
  p_e - p_\mu = \tfrac32 - 1 = \tfrac12, \qquad
  p_\mu - p_\tau = 1 - \tfrac23 = \tfrac13.
\]
These steps $\tfrac12, \tfrac13$ are the core: they generate the
$\sqrt{m}$-structure on which the Koide formula rests, and they are what
recurs throughout the following sections.

\section{Origin of the Prefactors: Resonances on the Torus}

\textbf{[K]} The prefactors $r_i$ are not freely chosen. Each lepton
corresponds to a standing field mode on the torus, characterised by two
integer winding numbers $(n_\theta, n_\phi)$ -- the number of circuits along
each cycle. This is the resonance condition: only modes with integer winding
are permitted, as with any standing wave on a ring.

\textbf{[K]} The prefactor follows directly from the winding numbers:
\[
  r_i = \frac{n_{\phi,i}^2}{n_{\theta,i}}.
\]
For the three leptons:
\[
  \begin{array}{lcccl}
    \text{Lepton} & n_\theta & n_\phi & r_i = n_\phi^2/n_\theta & \\[2pt]
    \text{Electron} & 3 & 2 & 4/3  & \\
    \text{Muon}     & 5 & 4 & 16/5 & \\
    \text{Tau}      & 9 & 5 & 25/9 &
  \end{array}
\]

\textbf{[K]} This explains the ``structure in the numerators'' visible in the $r_i$:
the numerators $4, 16, 25$ are the squares $n_\phi^2 = 2^2, 4^2, 5^2$,
the denominators $3, 5, 9$ are the $n_\theta$. The $r_i$ are simple
ratios of two resonance quantum numbers, not fitted values. The masses
of the sector thus follow from $\xi$, the ladder exponents, and a pair
of integer winding numbers per lepton.

\textbf{[K]} The same $r_i$ carry a second, equivalent interpretation: they are
the consonant intervals of just intonation. In just (5-limit) tuning,
consonant intervals are ratios of small integers with prime factors $2, 3, 5$
-- and that is exactly what the $r_i$ are:
\begin{center}
\begin{tabular}{lccl}
\hline
Lepton & $r_i$ & Prime factors & Just interval \\
\hline
Electron & $4/3$  & $2^2/3$   & Perfect fourth (498 ct) \\
Muon     & $16/5$ & $2^4/5$   & Min.\ sixth + octave (814 ct) \\
Tau      & $25/9$ & $5^2/3^2$ & Square of the major sixth \\
\hline
\end{tabular}
\end{center}

\textbf{[K]} This is not an analogy but the same fact seen from two sides:
a standing mode on a torus cycle \emph{is} a harmonic resonance,
like the overtones of a vibrating string. The winding numbers are
the order numbers of these resonances; the $r_i = n_\phi^2/n_\theta$ are the
corresponding interval ratios.

\textbf{[K]} The three charged leptons are the \emph{simplest} of these
resonances -- the low 5-limit consonances. The series does not stop there:
higher modes are the remaining fermions. The quarks are further overtones of the same torus:
the top is the only overtone above the Higgs ground tone, the bottom with $r=3/2$
-- the perfect fifth -- the harmonically simplest quark. There is therefore no
cutoff at three modes; there is an ascending resonance ladder, and the leptons
are its lowest part.

\textbf{[K]} The \emph{ordering} is therefore necessary, not chosen: resonances
are ordered by their order number, ascending winding number means
ascending mass. In a theory where particles are resonances,
``lightest resonance $=$ electron, next $=$ muon, then tau'' is not
an additional rule but the resonance structure itself.

\textbf{[K]} The tau prefactor $r_\tau = 25/9 = (5/3)^2$ is the square of the
major sixth. An earlier iteration carried $r_\tau = 8/3$ (giving
$m_\tau\approx1712$ MeV, $-3{.}6\,\%$); the correct value $25/9$ gives
$m_\tau\approx1783$ MeV and is recorded corpus-wide.

\section{The Three Mass Ratios from $\xi$}

\textbf{[K]} Because the masses are products $r_i\,\xi^{p_i}\,v$, the
absolute $v$ cancels in every ratio, leaving pure $\xi$-geometry:
\begin{align*}
  \frac{m_e}{m_\mu} &= \frac{r_e}{r_\mu}\,\xi^{p_e-p_\mu}
    = \frac{4/3}{16/5}\,\xi^{1/2} = \frac{5}{12}\,\xi^{1/2} = 0{.}004811, \\
  \frac{m_\mu}{m_\tau} &= \frac{r_\mu}{r_\tau}\,\xi^{p_\mu-p_\tau}
    = \frac{16/5}{25/9}\,\xi^{1/3} = \frac{144}{125}\,\xi^{1/3} = 0{.}05885, \\
  \frac{m_e}{m_\tau} &= \frac{r_e}{r_\tau}\,\xi^{p_e-p_\tau}
    = \frac{12}{25}\,\xi^{5/6} = 0{.}000283.
\end{align*}

\textbf{[K]} Against experiment (PDG):
\[
  \begin{array}{lccc}
    \text{Ratio} & \text{from } \xi & \text{PDG} & \text{dev.} \\[2pt]
    m_e/m_\mu   & 0{.}004811 & 0{.}004836 & 0{.}51\,\% \\
    m_\mu/m_\tau & 0{.}05885 & 0{.}05946 & 1{.}02\,\% \\
    m_e/m_\tau  & 0{.}000283 & 0{.}0002876 & 1{.}6\,\%
  \end{array}
\]
All three to roughly one percent, from a single parameter and rational numbers.

\section{Koide as Consequence: The Independent Route}

\textbf{[K]} The Koide relation
\[
  Q = \frac{m_e+m_\mu+m_\tau}{(\sqrt{m_e}+\sqrt{m_\mu}+\sqrt{m_\tau})^2}
\]
is not an independent empirical symmetry but follows from the ladder. Inserting
the three calculated masses -- \emph{without} using Koide anywhere --
gives
\[
  Q_\text{ladder} = 0{.}66774,
\]
which is $0{.}16\,\%$ from $2/3$. The geometric ladder thus predicts the
Koide quantity on its own. That the experimental value then hits exactly $2/3$
is an \emph{independent} confirmation of the exponent structure, not its
prerequisite.

\textbf{[B]} The reason $Q$ comes out so accurately is ratio-invariance: since
only exponent \emph{differences} enter and $v$ cancels, $Q$ depends solely on
$\xi^{p_i-p_j}$ and is free of any absolute scale. This is the purest form in
which $\xi$ appears.

\section{Three Routes to 2/3}

\textbf{[B]} At exactly $D=3$ (the topological dimension of the $T^4$ sector)
three different structures meet at the same value $2/3$:
\[
  1 - \frac1D = \frac23, \qquad \frac2D = \frac23, \qquad
  \frac12\Bigl(1+\frac1D\Bigr) = \frac23.
\]
The first is the complement step, the second the Koide equilibrium condition
between democratic and orthogonal projection components ($A=B$), the third
the ladder ratio. A formal deformation $D=3+\varepsilon$ separates the three again
-- they coincide only at integer $D=3$. This binds $2/3$ to the
\emph{integer} topological dimension, not to the fractal $D_f=3-\xi$.

\section{The Circulant: Sharpening to the Tau Mass}

\textbf{[K]} Alongside the ladder there is a second, complementary description:
the three $\sqrt{m_i}$ as discrete Fourier transforms of a
$Z_3$-circulant. In this representation two structural quantities appear,
an amplitude and a phase angle $\theta$, and the corpus finds
\[
  \theta = \frac29.
\]

\textbf{[K]} Two conditions then fix $m_\tau$: the Koide condition
$Q=2/3$ and the phase condition $\theta=2/9$. Both are formally distinct,
but at the current measurement precision they coincide (distance
$<0{.}05\,\sigma$), delivering the same $m_\tau$; the sharp prediction is
\[
  m_\tau^\text{pred} = 1776{.}968\ \text{MeV}.
\]

\textbf{[K]} The bookkeeping of the sector is decisive here. The lepton masses
are \emph{not} inputs to the theory: the structural numbers
($\sqrt2$ from $Q=2/3$, $2/9$ as address) are theory-side. To anchor
to the measured pole masses, exactly \emph{one} declared reference point is used:
the ratio $m_\mu/m_e$. It fixes the operative angle
$\theta_\text{ref} = 0{.}2222220471$ and is thereby consumed as a calibration quantity.

\textbf{[K]} The gain of this one anchor choice is the sharp prediction:
\[
  m_\tau = 1776{.}9690 \pm 0{.}0001\ \text{MeV},
\]
input-error-free, $+0{.}91\,\sigma$ against PDG $1776{.}86\pm0{.}12$
(Belle-II world average decides). This is the series' actual prediction
(A210, A220): a numerical value fixed \emph{before} measurement -- not a
bookkeeping value and not a scale choice like $\alpha$ or $G$, but a
genuinely measurable quantity, because the ratio structure forces the third mass
from the two light ones.

\textbf{[B]} The $Q$-route and the $\theta$-route are \emph{two} formally distinct
conditions, not one. At current measurement precision they coincide
(distance $<0{.}05\,\sigma$), delivering the same $m_\tau$; they are separated
only by a measurement that determines $m_\tau$ more sharply than this bound.

\section{The xi-Cascade: Why the Masses Are Most Precise}

\textbf{[B]} The lepton sector stands at the pure end of a hierarchy that orders
the whole theory:
\[
  \xi \xrightarrow{\text{ratios}} Q=\tfrac23
    \xrightarrow{\text{scale}} \alpha
    \xrightarrow{\text{SI metric}} G.
\]
Each step adds a layer of complexity, and precision falls accordingly:
\[
  \begin{array}{lcc}
    \text{Quantity} & \text{requires} & \text{precision} \\[2pt]
    Q\ (\text{Koide}) & \text{ratios only} & 0{.}00003\,\% \\
    \alpha & +\ \text{energy scale } E_0 & 0{.}006\,\% \\
    G & +\ \text{Planck length} & 0{.}5\,\%
  \end{array}
\]

\textbf{[K]} This explains the accuracy classes from A220 from \emph{one}
principle: the closer a quantity is to pure ratios, the more accurately it
follows from $\xi$; the more absolute scale it requires, the more bookkeeping
enters. The lepton masses are the purest form and therefore the most precise
and the only genuine prediction. $\alpha$ and $G$ are the same $\xi$-geometry,
only in increasingly dressed form (A130, A145).

\section{Verification Script}

\begin{itemize}
  \item The Yukawa ladder $m_i = r_i\xi^{p_i}v$ with the three mass ratios
    from $\xi$ (range check against PDG); the ladder prediction
    $Q_\text{ladder}=0{.}6677$ without Koide input; the two-route sharpening to
    $m_\tau$ (prediction, $0{.}90\,\sigma$); and the $\xi$-cascade with falling
    precision:
    \path{python/A_Serie_Skripte/a110_zirkulant.py}
\end{itemize}

\section{Symbol Glossary}

Symbols used throughout the A-Series are explained in A010.
Specific to this document:

\begin{center}
\renewcommand{\arraystretch}{1.3}
{\small\begin{tabular}{lp{9cm}}
\toprule
Symbol & Meaning \\
\midrule
$C$ & $3\times3$-circulant matrix; diagonalisable via DFT \\
$r=\sqrt{2}$ & eigenvalue ratio; condition for $Q=2/3$ \\
$Q = (\sum m_i)/(\sum\sqrt{m_i})^2$ & Koide operator; measured $Q\approx2/3$ \\
$\omega = e^{2\pi i/3}$ & primitive third root of unity \\
\bottomrule
\end{tabular}}
\renewcommand{\arraystretch}{1.0}
\end{center}

\section{Open Questions}

\textbf{First: the mode ordering is not an open point.} The three charged
leptons are the simplest torus resonances, and their ordering follows
necessarily from the ascending resonance order (above). Higher modes are the
remaining fermions; nothing cuts off at three. What is set -- and exclusively
that -- is the \emph{cardinality} three via the order $Z_3$ (A020). $Z_3$ is an
axiom, not a derivation, and an axiom needs no written derivation to be valid:
every theory enters its own framework at some point. $Z_3$ is motivated by the
parsimony of the three-point configuration in the two-dimensional cross-section
and by agreement with the generation number -- but these reasons do not make the
axiom a derivation, and their absence would not be a deficiency. What is to be
avoided is only the reversal: reading the three generations as a \emph{consequence}
of $Z_3$. In fact the observed number three motivates the choice $Z_3$, not vice versa.

\textbf{Second, the percent scatter of the ladder is not an open point but
the reason for the reference-point booking.} The ladder is a percent-level
statement; a $\xi$ back-calculated from it scatters accordingly. That is exactly
why the sector anchors the scale not in the ladder but in \emph{one}
reference point $m_\mu/m_e$. The circulant with the pair
$(\sqrt2, 2/9)$ misses $m_\mu/m_e$ at pole-mass level by $1{.}0\cdot10^{-5}$
-- which at this level excludes reading the pair as \emph{exact}; as a
rational address to seven digits it remains valid. The booking is thus
resolved, not open.

\textbf{Third, $\theta=2/9$ is precisely located, not unmotivated.} $2/9$
is the rational address of the reference angle $\theta_\text{ref}$ to seven
digits ($|\theta_\text{ref}-2/9| = 1{.}75\cdot10^{-7}$). It lives in the angular
phase sector and corresponds to the phase position $\chi=\pi/2$; the radial
mass sector remains rigid. What stands as a candidate -- not an achieved result
-- is only the \emph{dynamic selection}: whether a skew-adjoint generator outside
the static geometry locks the phase onto $\chi=\pi/2$. FFGFT's own real recursion
carries no such generator; whether an external dynamics supplies it is the only
remaining open point -- the location and address of $2/9$ are fixed.

\textbf{Fourth, the dependence on $m_\tau$ is the falsifiability, not a
deficiency.} After the reference-point booking, $m_\mu/m_e$ is the consumed
anchor and $m_\tau = 1776{.}9690\pm0{.}0001$ the one sharp test quantity. That
the sector stands or falls by it is intentional: a measurement placing $m_\tau$
clearly away from $1776{.}97$ refutes the exponent structure (A210). That is the
purpose of a prediction, not its weakness.

\section{Supersedes}

This document subsumes: Doc.~116 (Koide from $\xi$, Yukawa ladder, $\xi$-cascade),
Doc.~258 (three routes to $2/3$, $D=3+\varepsilon$), Doc.~259 (cross-terms),
Doc.~060 (musical resonance interpretation of prefactors), Doc.~189 (torus derivation of winding numbers), Doc.~190 (register K2 on $r_\tau$, R56 on $Z_3$), Doc.~282 (circulant, HLV bridge), Doc.~291 (dynamic location of $\theta=2/9$,
$\chi=\pi/2$), Doc.~292 (empirical lepton check).
Incorporated: P40 (ratios/bridge constants), P42 (one reference point $m_\mu/m_e$).

\section{Use of AI Tools}

This document was created with the assistance of AI tools. The questions,
stipulations, and all substantive decisions come from the author; the tools
formulate, compute, and create verification scripts.
Responsibility for content and errors rests entirely with the author.
Full wording of the rules in A010.
